\subsection{Directed dependency graph construction}
\label{sec:directed_dependency_graph}

The directed dependency graph is treated as a separate analytical object from the undirected
semantic-similarity graph. Its purpose is to estimate temporally admissible prerequisite relations,
not merely semantic proximity. The main design principle is:
\[
\textit{learn dependencies globally across trajectories, then export directed graphs locally.}
\]

\paragraph{Core rules.}
\begin{itemize}
    \item A skill may point only to a skill that appears later in the same trajectory.
    \item A skill in year $1$ cannot depend on a skill in year $2$ of the same trajectory.
    \item Skills taught in the same class and same year are \emph{not} directed unless a validated
    within-year order variable exists.
    \item The primary substantive graph is the local trajectory graph; the pooled graph is a
    higher-level summary only.
\end{itemize}

\begin{enumerate}

    \item \textbf{Step D1: Define the trajectory units and time order.}

    Construct a trajectory panel from the harmonized occurrence-level dataset. Each occurrence
    must be assigned to one ordered curricular path, denoted \texttt{trajectory\_unit}. The default
    rule is:
    \[
    \texttt{trajectory\_unit} =
    \begin{cases}
    \texttt{programme} \times \texttt{section}, & \text{if section defines a distinct path},\\
    \texttt{programme}, & \text{otherwise.}
    \end{cases}
    \]
    For each occurrence, retain:
    \texttt{trajectory\_unit}, \texttt{year}, \texttt{class\_id} (if available),
    \texttt{harmonized\_skill\_id}, and any possible within-year position variable.

    \textit{Expected outcome.}
    Every occurrence is assigned to exactly one valid trajectory and one valid temporal position.

    \textit{Files.}
    \texttt{10\_directed\_trajectory\_panel.csv},
    \texttt{10\_within\_year\_order\_audit.csv}.

    \item \textbf{Step D2: Create time-layered nodes.}

    Build the node table for the directed graph using time-layered nodes:
    \[
    \texttt{node} = (\texttt{trajectory\_unit}, \texttt{year}, \texttt{harmonized\_skill\_id}).
    \]
    If a validated within-year order exists, refine the node definition to:
    \[
    \texttt{node} =
    (\texttt{trajectory\_unit}, \texttt{year}, \texttt{within\_year\_position},
    \texttt{harmonized\_skill\_id}).
    \]
    This prevents the same skill from being collapsed across different years.

    \textit{Expected outcome.}
    The same skill may appear as different nodes in different years, preserving the temporal
    structure needed for a directed graph.

    \textit{Files.}
    \texttt{11\_directed\_nodes.csv},
    \texttt{11\_directed\_node\_provenance.csv}.

    \item \textbf{Step D3: Separate prerequisite candidates from same-time pairs.}

    For every pair of nodes within the same trajectory, classify the temporal relation as one of:
    \begin{itemize}
        \item earlier $\rightarrow$ later;
        \item same year / same class;
        \item same year / different class;
        \item reverse-time impossible relation.
    \end{itemize}
    Only earlier-later pairs remain eligible for the prerequisite graph. Same-time pairs are routed
    to an optional co-requisite layer. Reverse-time pairs are excluded.

    \textit{Expected outcome.}
    Only temporally admissible candidate prerequisite pairs remain in the directed branch.
    Same-time relations are preserved separately rather than forced into the DAG.

    \textit{Files.}
    \texttt{12\_temporal\_pair\_registry.csv},
    \texttt{12\_corequisite\_candidates.csv}.

    \item \textbf{Step D4: Compute temporal precedence evidence across trajectories.}

    For each skill pair $(i,j)$, compute the first appearance year in each trajectory:
    \[
    T_u(s) = \min \{ y : s \text{ appears in trajectory } u \text{ in year } y \}.
    \]
    Then compute:
    \[
    n^+_{ij} = \#\{u : T_u(i) < T_u(j)\}, \qquad
    n^-_{ij} = \#\{u : T_u(i) > T_u(j)\}, \qquad
    n^0_{ij} = \#\{u : T_u(i) = T_u(j)\}.
    \]
    Also record the precedence score:
    \[
    \mathrm{Prec}_{ij}
    =
    \frac{n^+_{ij} - n^-_{ij}}{n^+_{ij} + n^-_{ij} + n^0_{ij}}.
    \]

    \textit{Expected outcome.}
    A global table quantifies whether skill $i$ tends to appear before skill $j$, after it,
    or at the same time.

    \textit{Files.}
    \texttt{13\_precedence\_statistics.csv}.

    \item \textbf{Step D5: Build the direct-dependency learning sample.}

    For each child skill $j$, create a discrete-time learning sample based on first introduction.
    For each trajectory $u$ and year $y$, define:
    \[
    Y_{u,y,j}
    =
    \mathbf{1}\{\text{skill } j \text{ first appears in year } y+1\},
    \]
    and for each candidate parent skill $i$:
    \[
    X_{u,y,i}
    =
    \mathbf{1}\{\text{skill } i \text{ has appeared by year } y\}.
    \]
    Candidate parents are skills that are temporally admissible for $j$, meaning that they may
    appear before $j$ in at least one trajectory.

    \textit{Expected outcome.}
    Each child skill receives a tractable prediction dataset for estimating its direct parents.

    \textit{Files.}
    \texttt{14\_dependency\_learning\_sample.csv},
    \texttt{14\_candidate\_parent\_sets.csv}.

    \item \textbf{Step D6: Estimate direct parents with sparse conditional models.}

    For each child skill $j$, fit a sparse conditional model predicting first appearance from
    previously observed skills. A practical choice is a sparse logistic or discrete-time hazard model.
    The model is estimated only with predictors that satisfy the temporal constraint. Then repeat
    the estimation across bootstrap resamples of trajectory units and record the selection frequency
    of each admissible edge:
    \[
    q_{ij} = \Pr(i \rightarrow j \text{ is selected across bootstrap refits}).
    \]

    \textit{Expected outcome.}
    Each admissible edge $i \rightarrow j$ receives an empirical support score rather than a
    hand-made decision.

    \textit{Files.}
    \texttt{15\_direct\_parent\_scores.csv},
    \texttt{15\_bootstrap\_selection\_summary.csv}.

    \item \textbf{Step D7: Choose the edge cutoff automatically.}

    Convert the edge-support scores into a graph using an automatic cutoff rule. For each candidate
    threshold $\tau$ applied to $q_{ij}$, evaluate held-out prediction of first-introduction events.
    Select:
    \[
    \tau^* = \arg\max_{\tau} \mathrm{CVSkill}(\tau).
    \]
    If no thresholded graph outperforms the null model, do not force a binary graph; retain the
    weighted directed graph with edge weights $q_{ij}$.

    \textit{Expected outcome.}
    The graph cutoff is chosen by predictive performance rather than by heuristic topology rules.

    \textit{Files.}
    \texttt{16\_directed\_cutoff\_selection.csv},
    \texttt{16\_directed\_model\_selection\_surface.csv}.

    \item \textbf{Step D8: Construct local directed graphs.}

    For each trajectory unit $u$, instantiate an edge $i \rightarrow j$ when:
    \begin{enumerate}
        \item $q_{ij} \ge \tau^*$;
        \item both nodes are present in trajectory $u$;
        \item the temporal order in $u$ is admissible.
    \end{enumerate}
    Then apply transitive reduction within each local DAG so that redundant edges are removed.
    Export same-time relations separately as co-requisite edges.

    \textit{Expected outcome.}
    One directed dependency graph is produced per trajectory unit, plus an optional same-time
    co-requisite layer.

    \textit{Files.}
    \texttt{17\_directed\_edges\_trajectories.csv},
    \texttt{17\_corequisite\_edges\_trajectories.csv},
    \texttt{17\_directed\_graph\_summary.csv}.

    \item \textbf{Step D9: Aggregate upward for comparison.}

    Aggregate the local directed graphs to section-, track-, programme-, and pooled-level
    meta-graphs. For each higher-level edge $i \rightarrow j$, report summary support measures such
    as:
    \begin{itemize}
        \item frequency across local trajectories;
        \item mean edge-support probability $\bar q_{ij}$;
        \item inclusion rate across local units.
    \end{itemize}
    These higher-level graphs are descriptive summaries derived from the same learned local rule.

    \textit{Expected outcome.}
    Higher-level directed graphs become comparable because they all inherit the same local
    dependency estimator.

    \textit{Files.}
    \texttt{18\_directed\_edges\_sections.csv},
    \texttt{18\_directed\_edges\_tracks.csv},
    \texttt{18\_directed\_edges\_programmes.csv},
    \texttt{18\_directed\_edges\_pooled.csv},
    \texttt{18\_directed\_aggregation\_summary.csv}.

\end{enumerate}

\paragraph{Interpretation of same-year cases.}
If no validated within-year ordering variable exists, skills taught in the same class and same
year should not be forced into the prerequisite DAG. They belong in the co-requisite layer.
Only if a reliable within-year position variable exists should same-year directed edges be allowed.

\paragraph{Primary graph object.}
The primary substantive object is the local trajectory DAG. The section, track, programme, and
pooled graphs are aggregate summaries derived from the local DAG family. They should not be
treated as the only graph object.

\paragraph{Acyclicity note.}
Local trajectory graphs are DAGs by construction. Higher-level aggregated graphs may contain
directed cycles because they summarize heterogeneous local trajectories. They should therefore
be interpreted as weighted directional support graphs rather than as strict prerequisite DAGs.

\paragraph{Minimal expected outputs of the full directed branch.}
At the end of the directed-graph branch, the workflow should produce:
\begin{itemize}
    \item one trajectory panel with ordered occurrences;
    \item one time-layered node table;
    \item one temporal precedence table;
    \item one direct-parent score table with edge-support probabilities;
    \item one automatic cutoff-selection table;
    \item one directed graph per trajectory unit;
    \item one optional co-requisite layer for same-time relations;
    \item aggregated section, track, programme, and pooled directed meta-graphs.
\end{itemize}